\subsection{Componentes principales}\label{sec:pca}

Los predictores diarios no son coordenadas independientes: \texttt{NO2} y
\texttt{NOX} comparten la mayor parte de su variación, y otro tanto ocurre entre
\texttt{PM10} y \texttt{PM2.5}. El análisis de componentes principales no
persigue aquí un objetivo propio, sino sustituir ese bloque correlacionado por
un conjunto de ejes ortogonales que el modelo de clasificación pueda emplear sin
inflar los errores estándar de sus coeficientes.

\subsubsection{Bloque de predictores}

Entran las variables continuas de precursores y meteorología. Las indicadoras de
zona, \texttt{tipo\_dia}, \texttt{temporada} y \texttt{festivo} quedan fuera,
porque una columna binaria eleva de forma artificial la varianza explicada; se
incorporan directamente a la etapa de modelado. También se excluyen \ce{MDA8},
\texttt{O3\_max\_1h} y las indicadoras de excedencia, que constituyen la
respuesta y cuya inclusión como predictores sería fuga de información.

\texttt{PM2.5} se retira del bloque. La estación \texttt{NE3} no registra un
solo valor de esa variable y \texttt{NO3} registra 12.5\,\%, de modo que exigir
observaciones completas no descarta días aislados sino estaciones enteras: la
retención de la zona Noroeste caería de 60.5\,\% a 44.0\,\%, justamente la zona
de mayor tasa de excedencia. La dimensión de aerosoles permanece representada
por \texttt{PM10}. El bloque final reúne 15 predictores.

\subsubsection{Ajuste y criterio de retención}

La descomposición se ajusta sobre la matriz de correlación de los años 2021 a
2024 y se aplica al año 2025 mediante los mismos vectores propios y momentos de
entrenamiento; ajustarla sobre el periodo completo transferiría información del
conjunto de validación a la construcción de los ejes. El ajuste emplea 16{,}008
observaciones y la proyección de validación, 1{,}340.

El criterio de retención se fijó antes de ejecutar el análisis: método del codo
sobre el gráfico de sedimentación, presentado en \cref{fig:pca-scree}. La
curvatura del perfil de valores propios alcanza su máximo en la tercera
componente, con 1.13 frente a 0.26 en la cuarta, y las caídas sucesivas
confirman la lectura: el tramo pronunciado termina en la tercera componente y a
partir de ahí la pendiente se aplana. Se retienen tres componentes, que reúnen
55.9\,\% de la varianza.

El criterio de Kaiser retendría cinco componentes y 71.0\,\% de la varianza. Se
reporta como contraste, sin modificar la decisión: el umbral $\lambda > 1$ es
fijo y no considera la forma del perfil. Una lectura alternativa del codo, la de
máxima distancia a la recta que une los extremos, sitúa el punto en la cuarta
componente, si bien con un margen de 3\,\% sobre la tercera.

\begin{figure}[t]
  \centering
  \includegraphics[width=\linewidth]{pca_scree.pdf}
  \caption{Gráfico de sedimentación de los 15 predictores. El codo por
    curvatura se ubica en la tercera componente; la línea punteada marca el
    umbral de Kaiser.}
  \label{fig:pca-scree}
\end{figure}

\subsubsection{Interpretación de los ejes}

\cref{tab:pca-cargas} reúne las cargas dominantes de cada componente retenida.
Las tres admiten lectura física, lo que permite interpretar después los
coeficientes del modelo de clasificación en términos de condiciones
atmosféricas y no de componentes numeradas.

\begin{table}[h]
  \centering
  \caption{Cargas dominantes de las componentes retenidas, con la varianza que
    explica cada una.}
  \label{tab:pca-cargas}
  \begin{tabular}{lr}
    \toprule
    Componente y variable & Carga \\
    \midrule
    \multicolumn{2}{l}{\emph{Carga primaria de combustión} (26.1\,\%)} \\
    \quad \texttt{NOX}      & 0.88 \\
    \quad \texttt{NO}       & 0.82 \\
    \quad \texttt{NO2}      & 0.81 \\
    \quad \texttt{PM10}     & 0.66 \\
    \quad \texttt{WSR\_min} & $-0.48$ \\
    \midrule
    \multicolumn{2}{l}{\emph{Forzamiento fotoquímico} (19.8\,\%)} \\
    \quad \texttt{RH\_min\_diurna} & $-0.77$ \\
    \quad \texttt{TOUT\_max}       & 0.72 \\
    \quad \texttt{RH\_media}       & $-0.70$ \\
    \quad \texttt{SR\_acum}        & 0.65 \\
    \quad \texttt{WSR\_media}      & 0.60 \\
    \midrule
    \multicolumn{2}{l}{\emph{Estancamiento atmosférico} (10.0\,\%)} \\
    \quad \texttt{WSR\_min}        & $-0.50$ \\
    \quad \texttt{viento\_u\_media} & 0.48 \\
    \quad \texttt{WSR\_media}      & $-0.41$ \\
    \quad \texttt{PRS\_media}      & 0.39 \\
    \quad \texttt{RH\_min\_diurna} & $-0.37$ \\
    \bottomrule
  \end{tabular}
\end{table}

La primera componente agrupa los precursores de emisión vehicular en la ventana
matutina, acompañados de material particulado y de velocidad de viento con signo
negativo: describe la acumulación de carga primaria cuando el aire no la
dispersa. La segunda opone temperatura máxima y radiación acumulada a humedad
relativa, la firma de un día caluroso, despejado y seco. La tercera combina
viento débil con presión elevada, el patrón anticiclónico.

Este último eje merece una precisión metodológica. La combinación de viento
débil y presión alta no se construyó como índice para después presentarla como
resultado: emerge de la descomposición, que agrupa esas variables por su
covariación observada. \cref{fig:pca-biplot} muestra la estructura en los dos
planos definidos por la primera componente.

\begin{figure*}[t]
  \centering
  \includegraphics[width=\linewidth]{pca_biplot.pdf}
  \caption{Biplot de las componentes retenidas. Los puntos corresponden a una
    submuestra de 3{,}000 observaciones de entrenamiento, coloreadas por
    temporada; las flechas representan las cargas escaladas al marco del
    gráfico.}
  \label{fig:pca-biplot}
\end{figure*}

Los nombres describen la composición de cada eje y no atribuyen un mecanismo de
formación. El conjunto no incluye mediciones de compuestos orgánicos volátiles,
por lo que las componentes se describen como consistentes con procesos de
combustión, fotoquímica y estancamiento, sin sustentar afirmación alguna sobre
el régimen que limita la producción de ozono.

Queda una limitación del procedimiento. Los parámetros de la transformación
Box-Cox descrita en \cref{sec:transformaciones} se estimaron sobre el periodo
completo, de modo que la forma marginal de cada variable incorpora información
del año 2025. El efecto no alcanza a la variable respuesta ni a la construcción
de los ejes, que se ajustan únicamente con los años de entrenamiento.
